\documentclass[../../../include/open-logic-section]{subfiles}

\begin{document}

\olfileid[id]{sth}{ord-arithmetic}{add}
\olsection{Penjumlahan Ordinal}

Andaikan kita hendak menjumlahkan $\alpha$ dan~$\beta$. Kita cukup
menempatkan suatu salinan~$\beta$ tepat sesudah suatu salinan~$\alpha$.
(Kita perlu mengambil \emph{salinan}, sebab berdasarkan
\olref[ordinals][basic]{ordinalsaresubsets}, kita tahu bahwa
$\alpha\subseteq\beta$ atau $\beta\subseteq\alpha$.) Efek intuitifnya ialah
menjalani barisan langkah sepanjang~$\alpha$, \emph{lalu} menjalani barisan
sepanjang~$\beta$. Barisan yang dihasilkan akan terurut baik; jadi, berdasarkan
\olref[ordinals][ordtype]{thmOrdinalRepresentation}, barisan itu isomorfik
dengan suatu ordinal (unik). Ordinal tersebut dapat dipandang sebagai
\emph{jumlah} $\alpha$ dan~$\beta$.

Itulah gagasan intuitif di balik penjumlahan ordinal. Untuk mendefinisikannya
secara teliti, kita mulai dengan gagasan mengambil \emph{salinan} himpunan.
Gagasannya ialah menggunakan penanda sebarang, $0$ dan~$1$, untuk mencatat
asal setiap objek:

\begin{defn}\ollabel{defdissum}
\emph{Jumlah lepas} $A$ dan $B$ ialah
$A\disjointsum B=(A\times\{0\})\cup(B\times\{1\})$.
\end{defn}

Selanjutnya, kita mendefinisikan suatu pengurutan pada pasangan ordinal:

\begin{defn}
Untuk sebarang ordinal $\alpha_1,\alpha_2,\beta_1,\beta_2$, tetapkan bahwa:
\begin{align*}
	\tuple{\alpha_1, \alpha_2} \rlexless \tuple{\beta_1, \beta_2}\text{ jika dan hanya jika }& 
	\text{$\alpha_2 \in \beta_2$, atau}\\
	& \text{$\alpha_2 = \beta_2$ dan $\alpha_1 \in \beta_1$.}
\end{align*}
\end{defn}

Ini merupakan pengurutan \emph{leksikografis terbalik}, karena kita
mengurutkan menurut anggota kedua, lalu menurut anggota pertama. Ingat bahwa
kita hendak mendefinisikan $\alpha\ordplus\beta$ sebagai tipe urutan suatu
salinan~$\alpha$ yang diikuti suatu salinan~$\beta$. Untuk mencapainya, kita
menetapkan:

\begin{defn}\ollabel{defordplus}
Untuk sebarang ordinal $\alpha,\beta$, jumlahnya ialah
$\alpha\ordplus\beta=\ordtype{\alpha\disjointsum\beta,\rlexless}$.
\end{defn}
\noindent
Perhatikan bahwa kita sedikit menyalahgunakan notasi di sini; secara tegas,
kita seharusnya menulis
``$\Setabs{\tuple{x,y}\in\alpha\disjointsum\beta}{x\rlexless y}$'' sebagai
ganti ``$\rlexless$''. Namun, demi keringkasan, selanjutnya kita akan terus
menyalahgunakan notasi dengan cara ini.

Hasil berikut, bersama
\olref[ordinals][ordtype]{thmOrdinalRepresentation}, memastikan bahwa definisi
kita terbentuk dengan baik:

\begin{lem}\ollabel{ordsumlessiswo}
$\tuple{\alpha\disjointsum\beta,\rlexless}$ merupakan pengurutan baik bagi
sebarang ordinal $\alpha$ dan~$\beta$.
\end{lem}

\begin{proof}
Jelas bahwa $\rlexless$ terhubung pada $\alpha\disjointsum\beta$. Untuk
menunjukkan bahwa relasi itu berlandasan baik, tetapkan suatu himpunan tak
kosong $X\subseteq\alpha\disjointsum\beta$. Misalkan $Y$ himpunan bagian
dari~$X$ yang koordinat keduanya sekecil mungkin, yakni
$Y=\Setabs{\tuple{\gamma,i}\in X}{(\forall\tuple{\delta,j}\in X)i\leq j}$.
Sekarang pilih anggota~$Y$ dengan koordinat pertama terkecil.
\end{proof}
\noindent
Jadi, kita mempunyai definisi eksplisit yang baik bagi penjumlahan ordinal.
Berikut fakta yang tidak mengejutkan (ingat bahwa $1=\{0\}$ berdasarkan
\olref[z][infinity-again]{defnomega}):
\begin{prop}
$\alpha\ordplus1=\ordsucc{\alpha}$ bagi sebarang ordinal~$\alpha$.
\end{prop}

\begin{proof}
Perhatikan isomorfisme~$f$ dari
$\ordsucc{\alpha}=\alpha\cup\{\alpha\}$ ke
$\alpha\disjointsum1=(\alpha\times\{0\})\cup(\{0\}\times\{1\})$ yang
diberikan oleh $f(\gamma)=\tuple{\gamma,0}$ bagi $\gamma\in\alpha$, dan
$f(\alpha)=\tuple{0,1}$.
\end{proof}
\noindent
Selain itu, mudah ditunjukkan bahwa penjumlahan memenuhi syarat-syarat
rekursif tertentu:

\begin{lem}\ollabel{ordadditionrecursion}
Untuk sebarang ordinal $\alpha,\beta$, berlaku:
\begin{align*}
	\alpha\ordplus 0 &= \alpha\\
	\alpha \ordplus  (\beta\ordplus 1) &= (\alpha \ordplus  \beta) \ordplus  1\\
	\alpha  \ordplus  \beta &= \supstrict_{\delta < \beta}(\alpha \ordplus  \delta) && \text{jika $\beta$ ordinal limit.}
\end{align*}
\end{lem}

\begin{proof}
Kita memeriksa setiap kasus; pertama:
\begin{align*}
	\alpha \ordplus  0 
	& = \ordtype{(\alpha \times \{0\}) \cup (0 \times \{1\}), \rlexless} \\
	&= \ordtype{\alpha \times \{0\}, \rlexless}\\
	&= \alpha\\
	\alpha \ordplus (\beta \ordplus  1) 
	%&= \ordtype{(\alpha\times \{0\}) \cup (\ordtype{\beta \ordplus  1}\times \{1\}), \rlexless} \\
	&= \ordtype{(\alpha\times \{0\}) \cup (\ordsucc{\beta}\times \{1\}), \rlexless} \\
	&= \ordtype{(\alpha\times \{0\}) \cup (\beta \times \{1\}), \rlexless} \ordplus 1\\
	&= (\alpha \ordplus  \beta) \ordplus  1.
\end{align*}
Sekarang misalkan $\beta\neq\emptyset$ suatu limit. Jika $\delta<\beta$,
maka $\delta\ordplus1<\beta$, sehingga $\alpha\ordplus\delta$ merupakan
segmen awal sejati dari $\alpha\ordplus\beta$. Jadi,
$\alpha\ordplus\beta$ merupakan batas atas ketat bagi
$X=\Setabs{\alpha\ordplus\delta}{\delta<\beta}$. Selain itu, jika
$\alpha\leq\gamma<\alpha\ordplus\beta$, maka jelas
$\gamma=\alpha\ordplus\delta$ bagi suatu $\delta<\beta$. Jadi,
$\alpha\ordplus\beta=\supstrict_{\delta<\beta}(\alpha\ordplus\delta)$.
\end{proof}

Namun, berikut suatu fakta yang mencolok. Untuk mendefinisikan penjumlahan
ordinal, kita sebenarnya dapat menggunakan Teorema Rekursi Transfinit dan
menetapkan persamaan-persamaan rekursi tepat seperti dalam
\olref{ordadditionrecursion} (meskipun menggunakan ``$\ordsucc{\beta}$''
alih-alih ``$\beta\ordplus1$'').

Jadi, ada dua cara berbeda untuk mendefinisikan operasi pada ordinal. Kita
dapat mendefinisikannya \emph{secara sintetis}, dengan secara eksplisit
mengonstruksi suatu himpunan terurut baik dan mempertimbangkan tipe urutannya.
Atau, kita dapat mendefinisikannya \emph{secara rekursif}, hanya dengan
menetapkan persamaan-persamaan rekursi. Jika dilakukan dengan benar, hasilnya
sama. Teorema \olref[ordinals][ordtype]{thmOrdinalRepresentation} menjamin
bahwa persamaan-persamaan rekursi ini menentukan ordinal yang \emph{unik}.

Dalam banyak hal, aritmetika ordinal berperilaku seperti penjumlahan bilangan
asli. Sebagai contoh, kita dapat membuktikan hasil berikut:

\begin{lem}\ollabel{ordinaladditionisnice}
Jika $\alpha,\beta,\gamma$ ordinal, maka:
\begin{enumerate}
	\item\ollabel{ordaddition1} jika $\beta<\gamma$, maka
	$\alpha\ordplus\beta<\alpha\ordplus\gamma$;
	\item\ollabel{ordaddition2} jika
	$\alpha\ordplus\beta=\alpha\ordplus\gamma$, maka $\beta=\gamma$;
	\item\ollabel{ordaddition3}
	$\alpha\ordplus(\beta\ordplus\gamma)
	=(\alpha\ordplus\beta)\ordplus\gamma$, yakni penjumlahan bersifat
	asosiatif;
	\item\ollabel{ordaddition4} jika $\alpha\leq\beta$, maka
	$\alpha\ordplus\gamma\leq\beta\ordplus\gamma$.
\end{enumerate}
\end{lem}

\begin{proof}
Kita membuktikan \olref{ordaddition3} dan menyerahkan sisanya sebagai
latihan. Bukti dilakukan dengan Induksi Transfinit Sederhana pada~$\gamma$,
menggunakan \olref{ordadditionrecursion}. Jika $\gamma=0$:
\[
(\alpha \ordplus  \beta) \ordplus  0 = \alpha \ordplus  \beta  = \alpha \ordplus  (\beta \ordplus  0)
\]
Jika $\gamma=\delta\ordplus1$, andaikan sebagai hipotesis induksi bahwa
$(\alpha\ordplus\beta)\ordplus\delta
=\alpha\ordplus(\beta\ordplus\delta)$; sekarang gunakan
\olref{ordadditionrecursion} tiga kali:
\begin{align*}
	(\alpha \ordplus  \beta) \ordplus  (\delta \ordplus  1) & = ((\alpha \ordplus  \beta) \ordplus  \delta)\ordplus 1\\
	& = (\alpha \ordplus  (\beta \ordplus  \delta)) \ordplus  1\\
	& = \alpha \ordplus  ((\beta \ordplus  \delta)\ordplus 1)\\
	& = \alpha \ordplus  (\beta \ordplus  (\delta\ordplus 1)).
\end{align*}
Jika $\gamma$ ordinal limit, andaikan sebagai hipotesis induksi bahwa jika
$\delta\in\gamma$, maka
$(\alpha\ordplus\beta)\ordplus\delta
=\alpha\ordplus(\beta\ordplus\delta)$; sekarang:
\begin{align*}
	(\alpha \ordplus  \beta) \ordplus  \gamma & = \supstrict_{\delta < \gamma}((\alpha \ordplus  \beta) \ordplus  \delta) \\
	&= \supstrict_{\delta < \gamma}(\alpha \ordplus  (\beta \ordplus  \delta))\\
	&= \alpha \ordplus  \supstrict_{\delta < \gamma}(\beta \ordplus  \delta)\\
	& = \alpha \ordplus  (\beta \ordplus  \gamma).
\end{align*}
\end{proof}

\begin{prob}
Buktikan bagian selebihnya dari
\olref[sth][ord-arithmetic][add]{ordinaladditionisnice}.
\end{prob}

Dalam berbagai hal ini, penjumlahan ordinal seharusnya terasa sangat akrab.
Namun, ada satu perbedaan penting: penjumlahan ordinal \emph{tidak} seperti
penjumlahan bilangan asli.

\begin{prop}\ollabel{ordsumnotcommute}
Penjumlahan ordinal \emph{tidak} komutatif;
$1\ordplus\omega=\omega<\omega\ordplus1$.
\end{prop}

\begin{proof}
Perhatikan bahwa $1\ordplus\omega=\supstrict_{n<\omega}(1\ordplus n)
=\omega\in\omega\cup\{\omega\}=\ordsucc{\omega}=\omega\ordplus1$.
\end{proof}
\noindent
Walaupun pada mulanya hal ini mungkin mengejutkan, \emph{seharusnya tidak}.
Di satu sisi, ketika mempertimbangkan $1\ordplus\omega$, kita memikirkan tipe
urutan yang diperoleh dengan menempatkan satu anggota tambahan \emph{sebelum}
semua bilangan asli. Dengan penalaran seperti pada Hotel Hilbert dalam
\olref[sfr][infinite][hilbert]{sec}, secara intuitif anggota pertama tambahan
ini seharusnya tidak mengubah tipe urutan keseluruhan. Di sisi lain, ketika
mempertimbangkan $\omega\ordplus1$, kita memikirkan tipe urutan yang diperoleh
dengan menempatkan satu anggota tambahan \emph{sesudah} semua bilangan asli.
Itu merupakan sesuatu yang sama sekali berbeda!

\end{document}
